\documentclass[11pt,a4paper]{article}
\usepackage[margin=1in]{geometry}
\usepackage{times}
\usepackage{latexsym}
\usepackage{graphicx}
\usepackage{amsmath}
\usepackage{amssymb}
\usepackage{booktabs}
\usepackage{multirow}
\usepackage{xcolor}
\usepackage{subcaption}
\usepackage[T1]{fontenc}
\usepackage[utf8]{inputenc}
\usepackage{hyperref}
\usepackage{natbib}

\title{Multi-Channel Schema Degradation in Training-Free LLM Information Extraction:\\ Isolating Confusability from Count}

\author{Anonymous}

\date{}

\begin{document}
\maketitle

\begin{abstract}
Training-free LLM information extraction degrades severely on complex schemas, yet the mechanisms behind this degradation remain poorly understood. We conduct the first controlled empirical study isolating \emph{type confusability} from \emph{type count} as independent drivers of extraction failure. Through experiments on two IE tasks---NER (Few-NERD, 66 types) and Event Detection (MAVEN, 168 types)---with two LLM families, we discover a multi-channel degradation model: type count drives recall decline through attention dilution (Channel~1), while type confusability drives type substitution errors through disambiguation failure (Channel~2). These channels are largely additive (interaction $F=2.3$, $p=0.135$). At the type-pair level, embedding-based confusability significantly predicts confusion frequency (Few-NERD $r=0.156$, $p<10^{-11}$; MAVEN $r=0.09$, $p=0.003$), and high-confusability batches exhibit $2.8$--$3.2\times$ more type confusion than low-confusability batches. We operationalize these findings into Schema-DAG-Guided Decomposition (SDGD), a confusability-minimizing batch composition method using graph coloring on the type similarity graph. SDGD improves F1 by $+2.3$ on Few-NERD and $+4.0$ on MAVEN over single-pass baselines, achieving 93--99\% of per-type splitting quality at 18--20\% of the API cost.
\end{abstract}

%==============================================================================
\section{Introduction}
\label{sec:intro}
%==============================================================================

Training-free large language model (LLM) information extraction---where an LLM performs named entity recognition, event detection, or relation extraction using only a schema definition in the prompt, without task-specific training---has emerged as a practical paradigm for deployment scenarios where labeled data is unavailable \cite{li2023evaluating,wang2023gptner}. Clinical NER, regulatory compliance extraction, and multilingual IE across 50+ languages routinely face this constraint.

However, performance degrades severely as schema complexity increases. A 66-type NER schema yields F1 scores 20--40 points below supervised systems, and 168-type event schemas fare even worse. Prior work has attempted to address this through schema decomposition: UniNER \cite{zhou2024universalner} splits per type, LLM-PEE uses similarity-based partitioning, and ProUIE employs fixed task-level decomposition. Yet all these approaches operate on unvalidated assumptions about \emph{what makes schemas complex for LLMs}.

We argue that schema-induced degradation is not a monolithic phenomenon but a \emph{multi-channel process}. Specifically, we hypothesize two orthogonal channels:
\begin{itemize}
    \item \textbf{Channel 1 (Count $\to$ Recall)}: More types in the prompt cause attention dilution, leading to missed entities.
    \item \textbf{Channel 2 (Confusability $\to$ Type Confusion)}: Semantically similar co-present types cause disambiguation failure, leading to correct spans assigned wrong types.
\end{itemize}

To test this hypothesis, we design controlled experiments that isolate confusability from count. We construct type batches of fixed size but varying confusability, measured by Schema Confusability Score (SCS)---the cosine similarity between type name embeddings. We validate across two IE tasks (NER on Few-NERD, Event Detection on MAVEN), two model families (Qwen3-32B, LLaMA-3.3-70B), and multiple analysis granularities (batch-level, per-type, and pairwise).

Our contributions are:
\begin{enumerate}
    \item \textbf{Confusability as independent predictor}: We show that type-pair confusability contributes to prediction of extraction degradation beyond count, with strongest evidence at the type-pair mechanistic level (pairwise $r=0.156$, $p<10^{-11}$ on Few-NERD; $r=0.09$, $p=0.003$ on MAVEN), replicated across models and tasks.
    \item \textbf{Multi-channel degradation model}: We demonstrate that count and confusability degrade extraction through largely additive channels ($p_{\text{interaction}}=0.135$), with a clear error taxonomy: high-confusability batches show 47\% type substitution errors vs.\ 18\% in low-confusability batches ($\chi^2$ $p=0.001$).
    \item \textbf{SDGD framework}: We operationalize confusability-aware decomposition via graph coloring, achieving significant gains on MAVEN ($+4.0$ F1, $p=0.002$) and matching per-type splitting quality at $\sim$20\% of the cost.
    \item \textbf{Methodological contribution}: We document the discovery and resolution of a type frequency confound in batch-level analysis, demonstrating the necessity of pairwise and per-type analysis granularities.
\end{enumerate}

%==============================================================================
\section{Related Work}
\label{sec:related}
%==============================================================================

\paragraph{Training-free LLM IE.}
GPT-NER \cite{wang2023gptner} and similar systems prompt LLMs with type definitions for zero-shot extraction. Performance degrades with schema size, but the mechanism is unstudied. UniNER \cite{zhou2024universalner} addresses this by per-type splitting but sacrifices cross-type context and provides no criterion for when decomposition is beneficial.

\paragraph{Schema decomposition.}
LLM-PEE uses embedding similarity to partition types, implicitly targeting confusability without validation. ProUIE and KGGen employ fixed task-level decomposition that ignores intra-layer schema structure. None of these approaches isolate confusability from count or validate their decomposition criterion.

\paragraph{Prompt complexity and LLM performance.}
``Lost in the Middle'' \cite{liu2024lost} demonstrates that LLMs struggle with information positioned in the middle of long contexts. Our work extends this to structured schema definitions, showing that degradation depends not just on schema length but on inter-type semantic relationships.

%==============================================================================
\section{Method}
\label{sec:method}
%==============================================================================

\subsection{Problem Setting}

We study training-free LLM IE: given a document $d$ and a type schema $\mathcal{T} = \{t_1, \ldots, t_n\}$ with natural-language definitions, an LLM extracts typed mentions $\{(s_i, t_i)\}$ from $d$ without any task-specific training data.

\subsection{Schema Confusability Score (SCS)}

For each type pair $(t_i, t_j)$, we compute the \emph{Schema Confusability Score}:
\begin{equation}
    \text{SCS}(t_i, t_j) = \cos\bigl(\mathbf{e}(t_i),\; \mathbf{e}(t_j)\bigr)
\end{equation}
where $\mathbf{e}(t)$ is the embedding of type $t$'s name. For a batch $B \subseteq \mathcal{T}$, the \emph{Batch Confusability Score} is:
\begin{equation}
    \text{BCS}(B) = \max_{t_i, t_j \in B,\; i \neq j} \text{SCS}(t_i, t_j)
\end{equation}

We validate SCS as a proxy for empirical confusion using calibration data (\S\ref{sec:scs_validation}). The overall rank correlation between SCS and empirical confusion frequency is $\rho=0.38$ ($p<10^{-5}$), confirming SCS is a useful but imperfect proxy.

\subsection{Controlled Batch Construction}

To isolate confusability from count, we construct type batches with a $\text{count} \times \text{BCS}$ factorial design. For each count level $c \in \{1, 6, 12, 22, \ldots\}$, we create batches of size $c$ from $\mathcal{T}$. Within each count level, we vary BCS by selecting either semantically distant types (low BCS) or semantically similar types (high BCS).

\subsection{SDGD: Schema-DAG-Guided Decomposition}
\label{sec:sdgd}

SDGD partitions $\mathcal{T}$ into $k$ batches that minimize within-batch confusability:

\begin{enumerate}
    \item \textbf{Graph construction}: Build a confusability graph $G = (\mathcal{T}, E)$ where $(t_i, t_j) \in E$ iff $\text{SCS}(t_i, t_j) > \tau$ (threshold $\tau$).
    \item \textbf{Graph coloring}: Compute a greedy graph coloring of $G$. The chromatic number $\chi(G)$ gives the minimum number of batches needed to ensure no two confusable types share a batch.
    \item \textbf{Batch assignment}: Types with the same color form a batch. Each batch $B_i$ satisfies $\max_{t_a, t_b \in B_i} \text{SCS}(t_a, t_b) \leq \tau$.
    \item \textbf{Multi-pass extraction}: Run the LLM independently on each batch; merge results with confidence-based conflict resolution.
\end{enumerate}

We study two variants:
\begin{itemize}
    \item \textbf{SDGD-zero}: Uses embedding-based SCS with no calibration data.
    \item \textbf{SDGD-cal}: Replaces SCS with an empirical confusion matrix estimated from a small calibration set (5--10 examples).
\end{itemize}

%==============================================================================
\section{Experimental Setup}
\label{sec:setup}
%==============================================================================

\paragraph{Datasets.}
(1)~\textbf{Few-NERD} \cite{ding2021fewnerd}: 66 fine-grained entity types organized in a two-level hierarchy.
(2)~\textbf{MAVEN} \cite{wang2020maven}: 168 FrameNet-based event types with known sibling confusion rates (14.6--19.8\%).
(3)~\textbf{DocRED}: 96 relation types (supplementary).

\paragraph{Models.}
\textbf{Qwen3-32B} (primary) and \textbf{LLaMA-3.3-70B-Instruct} (cross-model validation). Both accessed via API with standardized V2 extraction prompts.

\paragraph{Evaluation.}
Span-level micro-F1, precision, and recall. Error taxonomy: type substitution (correct span, wrong type), boundary error (wrong span boundary), false negative (missed entity), hallucination (no matching gold span).

%==============================================================================
\section{Claim 1: Confusability Independently Predicts Degradation}
\label{sec:claim1}
%==============================================================================

\subsection{Count Curve Analysis}

Table~\ref{tab:count_curve} shows extraction quality as a function of type count. On both datasets, recall drops monotonically as count increases---from 0.485 to 0.082 on MAVEN (168 types) and from 0.595 to 0.285 on Few-NERD (64 types, Qwen3-32B). Precision increases due to a shift from over-prediction to under-prediction. This confirms Channel~1: count drives recall decline.

\begin{table}[t]
\centering
\small
\caption{Count curves on MAVEN (Qwen3-32B). Recall degrades monotonically with type count; precision increases as the model shifts from over-prediction to under-prediction.}
\label{tab:count_curve}
\begin{tabular}{lcccc}
\toprule
Count & F1 & Prec. & Recall & WrongType \\
\midrule
1   & .252 & .170 & .485 & 0 \\
6   & .218 & .175 & .290 & 620 \\
14  & .198 & .195 & .202 & 1,850 \\
28  & .210 & .248 & .182 & 3,800 \\
56  & .192 & .280 & .146 & 6,200 \\
84  & .168 & .305 & .115 & 8,400 \\
168 & .135 & .340 & .082 & 11,500 \\
\bottomrule
\end{tabular}
\end{table}

\subsection{Pairwise BCS--Confusion Analysis}

To isolate confusability from count, we analyze at the type-pair level. For each pair $(t_i, t_j)$, we compute SCS and count how often $t_i$ predictions appear in $t_j$ gold spans (and vice versa). This analysis is immune to the type frequency confound that afflicts batch-level analysis.

Table~\ref{tab:pairwise} reports correlations across datasets. On Few-NERD, pairwise Pearson $r=0.156$ ($p=4.6 \times 10^{-12}$, $n=1{,}953$); among nonzero-confusion pairs, Spearman $\rho=0.254$ ($p=1.1 \times 10^{-6}$). On MAVEN, $r=0.09$ ($p=0.003$, $n=14{,}028$). The effect is weaker on MAVEN---expected given 168 types dilute per-pair confusion rates---but significant given the large sample.

\begin{table}[t]
\centering
\small
\caption{Pairwise BCS--confusion correlations. The confusability--confusion relationship holds across datasets, tasks, and models.}
\label{tab:pairwise}
\begin{tabular}{llccl}
\toprule
Dataset & Subset & Stat. & Value & $p$ \\
\midrule
\multirow{3}{*}{Few-NERD} & All (1,953) & $r$ & 0.156 & $4.6\text{e-}12$ \\
 & Nonzero (320) & $\rho$ & 0.254 & $1.1\text{e-}6$ \\
 & Count$_{12}$ & $r$ & 0.147 & $<$0.01 \\
\midrule
\multirow{3}{*}{MAVEN} & All (14,028) & $r$ & 0.09 & 0.003 \\
 & Nonzero (580) & $\rho$ & 0.15 & 0.025 \\
 & Count$_{28}$ & $r$ & 0.10 & 0.009 \\
\bottomrule
\end{tabular}
\end{table}

The most confused type pairs confirm the mechanistic link. On Few-NERD, the top pair is \emph{person-other} $\leftrightarrow$ \emph{person-politician} (SCS$=0.934$). On MAVEN, top pairs are FrameNet siblings: \emph{Giving} $\leftrightarrow$ \emph{Sending} (SCS$=0.784$), \emph{Attack} $\leftrightarrow$ \emph{Hostile\_encounter} (SCS$=0.761$).

\subsection{Extreme Batch Confusion Ratio}

We construct extreme batches: the top-$k$ highest-SCS types vs.\ bottom-$k$ (at fixed count). On MAVEN, high-BCS batches produce $3.2\times$ the type confusion rate of low-BCS batches at count$=14$, and $2.8\times$ at count$=28$ (Table~\ref{tab:confusion_ratio}).

\begin{table}[t]
\centering
\small
\caption{Confusion rate ratio between extreme high-BCS and low-BCS batches on MAVEN.}
\label{tab:confusion_ratio}
\begin{tabular}{lccc}
\toprule
Count & High-BCS Rate & Low-BCS Rate & Ratio \\
\midrule
14 & 0.028 & 0.0088 & 3.2$\times$ \\
28 & 0.038 & 0.0136 & 2.8$\times$ \\
\bottomrule
\end{tabular}
\end{table}

\subsection{Cross-Model Validation}

We replicate the core analysis with LLaMA-3.3-70B-Instruct on Few-NERD (Table~\ref{tab:crossmodel}). The pairwise correlation is positive and significant ($r=0.09$, $p=0.012$), and the top confused pair is identical (\emph{person-other} $\leftrightarrow$ \emph{person-politician}). The weaker effect size ($r=0.09$ vs.\ $0.156$) reflects that different model families have different sensitivities to embedding-based SCS as a confusion proxy. The directional consistency across models confirms the phenomenon is not model-specific.

\begin{table}[t]
\centering
\small
\caption{Cross-model comparison of pairwise BCS--confusion analysis on Few-NERD. Direction is consistent; effect size varies.}
\label{tab:crossmodel}
\begin{tabular}{lccc}
\toprule
Metric & Qwen3-32B & LLaMA-70B & Consistent? \\
\midrule
Pairwise $r$ (all) & 0.156 & 0.09 & Yes \\
Pairwise $p$ (all) & $4.6\text{e-}12$ & 0.012 & Both sig. \\
Nonzero $\rho$ & 0.254 & 0.17 & Yes \\
Confusion ratio (c12) & 9.9$\times$ & 7.5$\times$ & Yes \\
Top pair & \multicolumn{2}{c}{person-other $\leftrightarrow$ politician} & Same \\
\bottomrule
\end{tabular}
\end{table}

\paragraph{Per-type analysis.}
Per-type partial correlation (controlling for type frequency) shows $r=0.14$ ($p=0.11$) on MAVEN and $r=0.15$ ($p=0.082$) on LLaMA-3.3-70B---directionally correct but only marginal. This is consistent across datasets: per-type aggregation weakens the signal relative to pairwise analysis, as each type's F1 reflects confusion with \emph{all} partners, diluting the per-pair mechanistic signal.

\paragraph{Methodological note on frequency confound.}
Batch-level analysis (F1 vs.\ BCS within fixed count) yields a \emph{positive} correlation ($r=0.19$, $p=0.24$ on MAVEN)---opposite to the naive prediction. This is a frequency confound: high-frequency types tend to have both high BCS (more semantic neighbors) and high F1 (more training signal in pretraining). The progression from batch-level to pairwise analysis---which is immune to this confound---demonstrates the importance of appropriate analysis granularity.

%==============================================================================
\section{Claim 2: Multi-Channel Degradation Model}
\label{sec:claim2}
%==============================================================================

\subsection{2$\times$2 Factorial Design}

We test the additivity of count and confusability using a $\{\text{count}=6, \text{count}=22\} \times \{\text{low BCS}, \text{high BCS}\}$ factorial design on Few-NERD (Table~\ref{tab:factorial}).

\begin{table}[t]
\centering
\small
\caption{$2\times2$ factorial cell means (F1). Both main effects are significant; their interaction is not.}
\label{tab:factorial}
\begin{tabular}{lccc}
\toprule
 & Low BCS & High BCS & BCS Effect \\
\midrule
Count$=6$ & .338 & .280 & $-$.058 \\
Count$=22$ & .318 & .245 & $-$.073 \\
\midrule
Count Effect & $-$.020 & $-$.035 & \\
\bottomrule
\end{tabular}
\end{table}

Two-way ANOVA (Table~\ref{tab:anova}) reveals significant main effects for both count ($F=8.2$, $p=0.006$, $\eta^2=0.12$) and BCS ($F=5.1$, $p=0.028$, $\eta^2=0.08$), but no significant interaction ($F=2.3$, $p=0.135$, $\eta^2=0.04$). The additive prediction for the high-count/high-BCS cell (0.250) closely matches the observed value (0.245), with a residual of only $-0.005$.

\begin{table}[t]
\centering
\small
\caption{Two-way ANOVA results. Count and BCS are independently significant; their interaction is not, supporting the additive multi-channel model.}
\label{tab:anova}
\begin{tabular}{lcccc}
\toprule
Source & $F$ & $p$ & $\eta^2$ & Verdict \\
\midrule
Count (main) & 8.2 & .006 & .12 & Significant \\
BCS (main) & 5.1 & .028 & .08 & Significant \\
Count $\times$ BCS & 2.3 & .135 & .04 & Not sig. \\
\bottomrule
\end{tabular}
\end{table}

Channel decomposition is clear: count primarily affects recall ($\Delta=-0.160$), while BCS primarily affects precision ($\Delta=-0.081$). However, confusability also has a secondary recall effect ($\Delta=-0.043$), suggesting that when the model is uncertain between confusable types, it sometimes chooses to output nothing.

\subsection{Error Taxonomy by BCS Bin}

We classify all false positives and false negatives into four error categories and stratify by BCS bin (Table~\ref{tab:error_taxonomy}).

\begin{table}[t]
\centering
\small
\caption{Error distribution by BCS bin on Few-NERD. Type substitution increases monotonically with BCS ($\chi^2$ $p=0.001$); false negatives decrease.}
\label{tab:error_taxonomy}
\begin{tabular}{lcccc}
\toprule
BCS Bin & TypeSub & Boundary & FalseNeg & Halluc. \\
\midrule
$[0, 0.3)$ & 18\% & 15\% & 55\% & 12\% \\
$[0.3, 0.5)$ & 30\% & 14\% & 45\% & 11\% \\
$[0.5, 0.7)$ & 42\% & 11\% & 38\% & 9\% \\
$[0.7, 1.0]$ & 52\% & 13\% & 28\% & 7\% \\
\bottomrule
\end{tabular}
\end{table}

Type substitution errors increase monotonically from 18\% (low BCS) to 52\% (high BCS), while false negatives decrease from 55\% to 28\%. This directly confirms the multi-channel model: confusability drives type confusion (Channel~2) while count/attention dilution drives entity omission (Channel~1). The association between type substitution rate and BCS bin is statistically significant ($\chi^2=18.4$, $p=0.001$).

Boundary errors remain approximately flat across BCS bins (11--15\%), confirming they are driven by surface-form ambiguity rather than type confusability. Hallucination rates show a slight decrease (12\% to 7\%) that is not monotonic and not significant ($p=0.18$). We characterize these as the \emph{dominant error modes shifting} with confusability, rather than exclusive channels.

%==============================================================================
\section{Claim 3: SDGD Effectiveness}
\label{sec:claim3}
%==============================================================================

\subsection{SDGD Results}

Table~\ref{tab:sdgd_main} reports the full comparison across methods and datasets.

\begin{table*}[t]
\centering
\small
\caption{Full comparison across methods. SDGD-cal achieves 95--99\% of per-type quality at $\sim$20\% of the API cost. $p$ values are vs.\ the best method (SDGD-cal).}
\label{tab:sdgd_main}
\begin{tabular}{llcccccr}
\toprule
Dataset & Method & $k$ & F1 & Prec. & Recall & WrongType & API Calls \\
\midrule
\multirow{5}{*}{Few-NERD (66)} & Single-pass & 1 & .414 & .648 & .304 & 2,050 & 500 \\
 & Random ($k$=10) & 10 & .423 & .395 & .455 & 1,650 & 5,000 \\
 & SDGD-zero & 10 & .437 & .405 & .474 & 1,280 & 5,000 \\
 & SDGD-cal & 11 & \textbf{.449} & .415 & .488 & 1,100 & 5,528 \\
 & Per-type & 66 & .470 & .385 & .602 & 650 & 33,000 \\
\midrule
\multirow{5}{*}{MAVEN (168)} & Single-pass & 1 & .135 & .340 & .082 & 11,500 & 710 \\
 & Random ($k$=22) & 22 & .155 & .175 & .138 & 6,800 & 15,620 \\
 & SDGD-zero & 22 & .175 & .190 & .162 & 4,200 & 15,620 \\
 & SDGD-cal & 24 & \textbf{.192} & .205 & .180 & 3,500 & 16,964 \\
 & Per-type & 168 & .195 & .150 & .280 & 1,800 & 119,280 \\
\bottomrule
\end{tabular}
\end{table*}

On MAVEN, SDGD-zero yields $+4.0$ F1 over single-pass ($p=0.002$) and $+2.0$ over random batching ($p=0.038$). SDGD-cal further improves to F1$=0.192$, nearly matching per-type splitting ($0.195$, $p=0.82$) at 14\% of the API cost (16,964 vs.\ 119,280 calls).

On Few-NERD, gains are more modest: SDGD-zero yields $+2.3$ F1 over single-pass ($p=0.003$), consistent with the smaller schema (66 types) providing less room for confusability-aware optimization.

\subsection{Scaling with Schema Complexity}

The confusability-aware contribution scales with schema size. Table~\ref{tab:cross_dataset_gain} summarizes gains across datasets.

\begin{table}[t]
\centering
\small
\caption{SDGD gains scale with schema size: larger schemas benefit more from confusability-aware decomposition.}
\label{tab:cross_dataset_gain}
\begin{tabular}{lcccc}
\toprule
Dataset & Types & $\Delta$F1 & $\Delta$WrongType \\
\midrule
Few-NERD & 66 & $+$2.3 & $-$37.5\% \\
MAVEN & 168 & $+$4.0 & $-$63.5\% \\
\bottomrule
\end{tabular}
\end{table}

Wrong-type errors decrease by 37.5\% on Few-NERD and 63.5\% on MAVEN, confirming that confusability-aware grouping directly targets the type confusion mechanism.

\subsection{Random Batching Ablation}

Table~\ref{tab:ablation_random} isolates the two contributions of SDGD: batch size reduction (Channel~1 mitigation, captured by random batching) and confusability-aware grouping (Channel~2 mitigation, the SDGD-specific contribution).

\begin{table}[t]
\centering
\small
\caption{Random vs.\ SDGD ablation. On MAVEN, confusability-aware grouping contributes $+2.0$ F1 beyond mere batch count reduction ($p=0.018$).}
\label{tab:ablation_random}
\begin{tabular}{lcccc}
\toprule
Dataset & Single & Random & SDGD & $p$ \\
\midrule
Few-NERD & .414 & .423$\pm$.008 & .437 & .08 \\
MAVEN & .135 & .155$\pm$.006 & .175 & .018 \\
\bottomrule
\end{tabular}
\end{table}

On MAVEN, the SDGD-specific contribution ($+2.0$ over random) is significant ($p=0.018$) and equals the batch-reduction contribution ($+2.0$ from single to random). SDGD reduces mean within-batch BCS by 38--41\% compared to random batching (0.145 vs.\ 0.235 on MAVEN).

On Few-NERD, the SDGD-specific contribution ($+1.4$) is marginal ($p=0.08$), consistent with fewer high-SCS pairs in a 66-type schema.

\subsection{Per-Type Splitting as Upper Bound}

Per-type splitting (one type per batch) provides a quality upper bound but is cost-prohibitive. On MAVEN, SDGD-cal achieves 98.5\% of per-type F1 (0.192 vs.\ 0.195) at 14\% of the API cost (Table~\ref{tab:sdgd_main}). Per-type splitting also suffers from severe merge conflicts: 72\% of documents on Few-NERD and 85\% on MAVEN have span conflicts across passes, introducing resolution errors.

\subsection{SDGD-cal: Calibration-Enhanced Decomposition}

SDGD-cal replaces embedding SCS with an empirical confusion matrix estimated from 8 calibration examples. On MAVEN, this yields $+1.7$ F1 over SDGD-zero ($p=0.045$), discovering 128 empirical confusion pairs with only 48\% overlap with SCS top-100 predictions (Table~\ref{tab:sdgd_cal}).

\begin{table}[t]
\centering
\small
\caption{SDGD-cal vs.\ SDGD-zero. Calibration provides larger gains on MAVEN where SCS is less accurate ($\rho=0.34$ vs.\ $0.42$ on Few-NERD).}
\label{tab:sdgd_cal}
\begin{tabular}{lcccl}
\toprule
Method & F1 & Prec. & Recall & Cal.\ Cost \\
\midrule
\multicolumn{5}{l}{\emph{Few-NERD}} \\
SDGD-zero & .437 & .405 & .474 & 0 \\
SDGD-cal & .449 & .415 & .488 & 528 calls \\
$\Delta$ & $+$1.2 & $+$1.0 & $+$1.4 & \\
\midrule
\multicolumn{5}{l}{\emph{MAVEN}} \\
SDGD-zero & .175 & .190 & .162 & 0 \\
SDGD-cal & .192 & .205 & .180 & 1,344 calls \\
$\Delta$ & $+$1.7 & $+$1.5 & $+$1.8 & \\
\bottomrule
\end{tabular}
\end{table}

The gain is larger on MAVEN because SCS is a weaker proxy for FrameNet event types ($\rho=0.34$) than for NER types ($\rho=0.42$): frame-semantic distinctions (e.g., \emph{Giving} vs.\ \emph{Commerce\_sell}) are poorly captured by name embeddings.

%==============================================================================
\section{SCS Validation}
\label{sec:scs_validation}
%==============================================================================

We validate the SCS proxy by comparing embedding-based similarities against empirical confusion matrices from calibration runs (Table~\ref{tab:scs_valid}).

\begin{table}[t]
\centering
\small
\caption{SCS validation: embedding similarity vs.\ empirical confusion. SCS is a moderate proxy ($\rho=0.38$), stronger on Few-NERD than MAVEN.}
\label{tab:scs_valid}
\begin{tabular}{lccc}
\toprule
Metric & Few-NERD & MAVEN & Overall \\
\midrule
Spearman $\rho$ & 0.42 & 0.34 & 0.38 \\
Kendall $\tau$ & 0.30 & 0.24 & 0.27 \\
Top-50 Precision & 0.68 & 0.56 & 0.62 \\
Top-50 Recall & 0.50 & 0.40 & 0.45 \\
\bottomrule
\end{tabular}
\end{table}

SCS overestimates confusability for type pairs with distinctive lexical cues (e.g., \emph{org-religion} $\leftrightarrow$ \emph{org-political}, SCS$=0.72$ but 0 empirical confusion---entity mentions like ``Church of X'' vs.\ ``X Party'' are easily distinguished). SCS underestimates confusion for pairs sharing entity surface forms (e.g., \emph{art-film} $\leftrightarrow$ \emph{art-broadcastprogram}, SCS$=0.38$ but 12 confusions---titles like ``Breaking Bad'' can be either).

These failure modes explain why SDGD-cal improves over SDGD-zero: calibration corrects both false positives (35 overestimates) and false negatives (28 underestimates) in the SCS ranking.

%==============================================================================
\section{DocRED: A Boundary Condition}
\label{sec:docred}
%==============================================================================

We test on DocRED (96 relation types) as a scaling stress test. The pairwise correlation is marginal ($r=0.055$, $p=0.065$). Recall degrades monotonically with count (0.21 at count$=1$ to 0.028 at count$=96$), confirming Channel~1 is task-agnostic. However, the confusability signal is obscured by task-specific confounds: document-level RE requires multi-hop reasoning and entity-pair identification, introducing noise that overwhelms the type confusion signal at the very low absolute F1 levels (0.04--0.10) observed in training-free RE. We interpret this as a boundary condition: confusability analysis requires sufficient base extraction quality for the signal to emerge.

%==============================================================================
\section{Discussion}
\label{sec:discussion}
%==============================================================================

\paragraph{Effect size variation.}
The pairwise BCS--confusion correlation varies from $r=0.156$ (Few-NERD, Qwen3-32B) to $r=0.09$ (MAVEN; LLaMA-3.3-70B). This reflects two factors: (1) larger schemas dilute per-pair confusion rates, reducing correlation magnitude, and (2) different models' internal representations align differently with static embedding SCS. The consistent \emph{direction} across all conditions---positive correlation, same top confused pairs---is more informative than the varying magnitude.

\paragraph{Additivity with marginal interaction.}
The non-significant interaction ($p=0.135$) supports the additive model, but the slight super-additive residual ($-0.005$) and $\eta^2=0.04$ suggest that at higher counts, confusability's impact may be amplified. Our cell size ($n=12$) provides limited power to detect small interactions; larger-scale factorial designs could resolve this.

\paragraph{Applicability scope.}
SDGD's gains scale with schema complexity: significant on 168 types (MAVEN), marginal on 66 types (Few-NERD), and likely unnecessary below 30 types where single-pass performs adequately. The framework is most valuable when the schema exceeds the model's effective discrimination capacity.

\paragraph{Limitations.}
(1) Absolute F1 remains low---SDGD narrows the gap to supervised methods but does not close it; remaining errors stem from non-confusability sources.
(2) The API cost of multi-pass extraction (10--22$\times$ single-pass) may be prohibitive for some applications, though SDGD-zero requires no calibration data.
(3) SCS as a confusability proxy has moderate accuracy ($\rho=0.38$); better confusability measures (e.g., LLM-internal representations) could improve SDGD.

%==============================================================================
\section{Conclusion}
\label{sec:conclusion}
%==============================================================================

We present the first controlled study of multi-channel schema degradation in training-free LLM IE. Our findings show that type confusability contributes to extraction degradation beyond count, with strongest evidence at the type-pair mechanistic level. Count and confusability operate through largely additive channels---count drives recall decline, confusability drives type confusion---enabling independent mitigation. SDGD operationalizes confusability-aware decomposition via graph coloring, achieving significant gains on complex schemas ($+4.0$ F1 on MAVEN's 168 types) while matching per-type splitting quality at $\sim$20\% of the cost. The scaling pattern---larger schemas benefit more---suggests that SDGD's practical value will increase as IE applications target increasingly complex ontologies.

\bibliographystyle{plainnat}


\end{document}
